Read in Gauge Theory Gravity, $T_{p}M$ below is the space of physical vectors at $p$ with the flat product, the derived metric there, and $\nabla_{a}$ on a scalar is the covariant derivative $D$ (Fact \claimref{synthesis.imported-gtg-dictionary}, clauses (i) and (ii)). \emph{(i).} Work in $T_{p}M$ with signature $(-,+,+,+)$, $Q(v) = v\cdot v$, and the reflection $R_{v}(x) = x - 2\,(x\cdot v)\,v/Q(v)$ in a vector $v$ with $Q(v) \neq 0$. For vectors $a$, $b$ with $Q(a) = Q(b)$ and $Q(a - b) \neq 0$, $R_{a - b}(a) = b$, since $Q(a - b) = 2(Q(a) - a\cdot b)$. Let $\ell$, $n$ be future null with $\ell\cdot n \neq 0$ (non-parallel); the local Rindler horizon of Fact \claimref{synthesis.imported-jacobson-clausius-1995} has edge element $\Pi = \{\ell, n\}^{\perp}$, with $n$ the second null normal of the edge. Step 1: if $\ell$ is a positive multiple of $\ell_{0}$ (equivalently $\ell_{0}\cdot\ell = 0$), the rays already agree and step 1 is skipped. Otherwise put $v_{1} = \ell_{0} - \ell$; then $Q(v_{1}) = -2\,\ell_{0}\cdot\ell > 0$ (two non-parallel future null vectors have negative product), so $v_{1}$ is spacelike and $R_{v_{1}}(\ell_{0}) = \ell$. Write $c\ell$ ($c > 0$, $c = 1$ after step 1) for the image of $\ell_{0}$. Step 2: let $n_{0}'$ be the image of $n_{0}$, future null and not parallel to $\ell$, and $\mu = (n_{0}'\cdot\ell)/(n\cdot\ell)$, positive because both products are negative. If $v_{2} = n_{0}' - \mu n \neq 0$, then $v_{2}\cdot\ell = 0$, so $v_{2}$, orthogonal to a null vector and not proportional to it, is spacelike, $R_{v_{2}}$ fixes $\ell$, and $R_{v_{2}}(n_{0}') = \mu n$ (equal squares, zero). The composite $\Lambda$ is a product of at most two reflections in spacelike vectors; a product of two such reflections lies in the proper orthochronous Lorentz group, and when only one is used it is composed with the reflection in $e_{2}$, which fixes $\ell_{0}$, $n_{0}$ and $\Pi_{0}$ as a set; so $\Lambda\ell_{0} = c\ell$, $\Lambda n_{0} = \mu n$, and $\Lambda$ carries $\Pi_{0} = \{\ell_{0}, n_{0}\}^{\perp}$ isometrically onto $\{\ell, n\}^{\perp} = \Pi$, with its area element. So the group acts transitively on (future null ray, edge plane). The boosts of the $(\ell_{0}, n_{0})$ plane are products of reflections in $e_{0}$ and in $c'e_{0} + s'e_{1}$, $c'^{2} - s'^{2} = 1$, which map $\ell_{0} \mapsto (c' + s')^{2}\ell_{0}$, $n_{0} \mapsto (c' - s')^{2}n_{0}$ and fix $e_{2}$, $e_{3}$; the scale of $\ell$ is the normalisation of the approximate boost field, through which the acceleration $a$ of the observers labels it, so the area element is unchanged under any change of $a$. Under (Loc) the density at $p$ is a function on the set of data modulo isometries, which by transitivity is one point; any two tangent spaces are isometric, so the constant is the same at every $p$. The script checks the construction exactly for all ordered non-parallel pairs from nine rational null vectors (Pythagorean quadruples and two positive multiples of $\ell_{0}$), including positivity of $c$ and $\mu$, orthonormality of the image of $\langle e_{2}, e_{3}\rangle$, its orthogonality to $\ell$ and $n$, and the time orientation, and checks the boost rescaling for three rational boosts.

\emph{(ii).} Under (P), (R) and Hawking's hypothesis, Theorem \claimref{synthesis.quarter-from-temperature} gives $S(M) = \int_{0}^{M}8\pi M'\,dM' = 4\pi M^{2} = A/4$ with $A = 16\pi M^{2}$, and, as that theorem records, $S = (p\kappa/2\pi)\,A/4$ for a period $p$ in place of $2\pi/\kappa$. By the identification hypothesis, $S(M) = \int s\,dA + o(A) = \eta A + o(A)$ (the cross-section of the horizon at each of its points has flat tangent data of the kind in (i), so (Loc) assigns it the density $\eta$), whence $\eta = \lim_{M\to\infty}S/A = \tfrac14$, or $p\kappa/8\pi$ for the period $p$; the additive constant of $S$ and all subleading terms drop out in the limit. That patches of a Killing horizon are local Rindler horizons at leading order is shown in Theorem \claimref{synthesis.c-loc-is-a-second-placement-of-the-count}; the identification hypothesis makes it unnecessary here. The script recomputes $\kappa = f'(2M)/2$ as a rational function of $M$ from $f = 1 - 2M/r$, derives $1/T = 2\pi/\kappa = 8\pi M$ from it, checks that $4\pi M^{2}$ has derivative $1/T$ and vanishes at $M = 0$, that $S/A = \tfrac14$ identically in $M$, and that a period $c$ times the regular one gives $c/4$. For de Sitter, Theorem \claimref{synthesis.de-sitter-entropy-pi-n-squared} gives $S_{\Lambda} = \pi N^{2}$ and $A_{\Lambda} = 4\pi N^{2}$, ratio $\tfrac14$; Fact \claimref{synthesis.imported-kerr-surface-gravity-entropy} gives $A_{+}/4$ at every spin (the script evaluates a Kerr case with rational $r_{+}$). Finally, with $\eta = 4\pi\nu\langle s\rangle$ as in Theorem \claimref{synthesis.clausius-route-newton-constant-is-residue}(ii), $\eta = \tfrac14$ if and only if $16\pi\nu\langle s\rangle = 1$, which by part (iii) of that theorem is the area-law condition defining the registered $\nu$.

\emph{(iii).} If clause (i) of version 1 of Premise \claimref{synthesis.c-loc-local-horizon-premise} holds, the density is the constant $\eta$, which depends on nothing, so (Loc) holds. Conversely (Loc) gives, by (i), a universal $\eta$ per unit area, and by (ii) $\eta = \tfrac14$: this is hypothesis (b) of Fact \claimref{synthesis.imported-jacobson-clausius-1995} with the value the area law needs. (Loc) says nothing about which degrees of freedom carry the entropy, so the statement that boundary modes at the epoch density do is not implied by it.

\emph{(iv).} Take a patch of a local Rindler horizon through $p$ with generators $k$ of affine parameter $\lambda \in [\lambda_{0}, 0]$, $\lambda = 0$ at $p$, as in Fact \claimref{synthesis.imported-jacobson-clausius-1995}. The expansion vanishes at $p$, so the area element along a generator is $dA(\lambda) = dA_{p}(1 + O(\lambda^{2}))$, and the heat $\delta Q = \int T_{ab}\chi^{a}k^{b}\,d\lambda\,dA$ with $\chi = -\kappa\lambda k + O(\lambda^{2})$ is $O(\lambda_{0}^{2})$ per unit area. With entropy $\int\eta\,dA$ and $\eta(\lambda) = \eta(p) + \lambda\,k^{a}\nabla_{a}\eta + O(\lambda^{2})$, the entropy change per unit area from $\lambda_{0}$ to $0$ is $\eta(p) - \eta(\lambda_{0})(1 + O(\lambda_{0}^{2})) = -\lambda_{0}\,k^{a}\nabla_{a}\eta + O(\lambda_{0}^{2})$. The Clausius relation $\delta Q = T\,\delta S$ with $T = \kappa/2\pi \neq 0$ at first order in $\lambda_{0}$ therefore requires $k^{a}\nabla_{a}\eta = 0$ at $p$ for every future null $k$. Future null vectors span $T_{p}M$, so $\nabla\eta = 0$ at every point, and $\eta$ is constant on a connected region. The script checks the order counting with polynomials in $\lambda_{0}$: the first-order coefficient of $T\,\delta S - \delta Q$ is $-T\,k^{a}\nabla_{a}\eta$, and four null vectors with independent components determine the gradient.

\emph{Remark (argued, not part of the theorem).} (a) The calibration in (ii) uses the Schwarzschild solution with the constant $G$ that defines the Planck cell, so on Jacobson's route it returns $G_{\mathrm{cl}} = 1/(4\eta) = G$, a consistency and not a determination of $G$, as in Theorem \claimref{synthesis.clausius-route-newton-constant-is-residue}. (b) The curvature and the matter at $p$ are point values that do not become small as the patch shrinks, so independence from them is an assumption of (Loc), not a consequence of working at leading order in the patch size. A density depending on them, such as the Wald entropy of a Lagrangian with Ricci-polynomial terms, is excluded by (Loc) by assumption; such dependence would enter as $\ell_{P}^{2}|\mathrm{Riem}|$ or $\ell_{P}^{4}|T|$ in Planck units and is small at sub-Planckian curvature, but the theorem does not use that estimate. A dependence on the acceleration, the null direction or the rotation of a hole, as in refutation condition (a) of Premise \claimref{synthesis.c-loc-local-horizon-premise}, is excluded by (Loc). (c) The identification hypothesis of (ii) takes the hole's whole thermodynamic entropy to be the integral of the density. If the cutoff modes of Premises \claimref{singularity.c1-curvature-cutoff} and \claimref{singularity.c2-cell-density} carried the area-law entropy independently of the horizon density, the total would be twice it; so (Loc) with the calibration excludes independent placement and is compatible only with identification (one mode set) or a correlated joint state, by the same area-law reasoning as the placement result, Theorem \claimref{synthesis.c-loc-is-a-second-placement-of-the-count}. Whether the cutoff count equals $A_{+}/4$ on Kerr, as identification requires, is not decided by (Loc).
